\documentclass[11pt, a4paper]{article}

% ── Packages ──────────────────────────────────────────────────────────
\usepackage[utf8]{inputenc}
\usepackage[T1]{fontenc}
\usepackage{lmodern}
\usepackage[margin=1in]{geometry}
\usepackage{amsmath, amssymb, amsthm}
\usepackage{booktabs}
\usepackage{longtable}
\usepackage{graphicx}
\usepackage{hyperref}
\usepackage{xcolor}
\usepackage{listings}
\usepackage{tcolorbox}
\usepackage{polya}

\hypersetup{
  colorlinks=true,
  linkcolor=polyablue,
  urlcolor=polyablue,
  citecolor=polyablue,
}

% ── Code listing style ────────────────────────────────────────────────
\lstdefinestyle{polya-python}{
  language=Python,
  basicstyle=\ttfamily\small,
  keywordstyle=\color{polyablue}\bfseries,
  commentstyle=\color{polyagray}\itshape,
  stringstyle=\color{polyagreen},
  numbers=left,
  numberstyle=\tiny\color{polyagray},
  numbersep=8pt,
  frame=single,
  framerule=0.4pt,
  rulecolor=\color{polyagray},
  breaklines=true,
  showstringspaces=false,
  tabsize=4,
}
\lstset{style=polya-python}
\title{SIR Epidemic Model}
\author{uber-polya}
\date{2026-02-26}

\begin{document}

\maketitle
\tableofcontents
\newpage

% ── Phase A: Problem Formulation ─────────────────────────────────────
\section{Problem Formulation}

\begin{polyamodel}

\textbf{Problem Type}: Find \hfill
\textbf{Domain}: Dynamical Systems / ODEs


\end{polyamodel}

% ── Universe ──────────────────────────────────────────────────────────
\subsection{Universe}

\begin{itemize}
  \item $Population N = 10000$
  \item $Time horizon: [0, 200.0] days$
\end{itemize}

% ── Variables ─────────────────────────────────────────────────────────
\subsection{Variables}

\begin{longtable}{lll}
\toprule
\textbf{Symbol} & \textbf{Meaning} & \textbf{Type / Range} \\
\midrule
\endhead
$S(t)$ & susceptible individuals at time t & $[0, 10000]$ \\
$I(t)$ & infected individuals at time t & $[0, 10000]$ \\
$R(t)$ & recovered individuals at time t & $[0, 10000]$ \\
$\beta$ & transmission rate = 0.3/day & $\mathbb{{R}}_{{>0}}$ \\
$\gamma$ & recovery rate = 0.1/day & $\mathbb{{R}}_{{>0}}$ \\
$R_0$ & basic reproduction number = 0.3/0.1 & $\mathbb{{R}}_{{>0}}$ \\
$N$ & total population = 10000 & $\mathbb{{Z}}^+$ \\
\bottomrule
\end{longtable}

% ── Structure ─────────────────────────────────────────────────────────
\subsection{Structure}

A disease outbreak in a population of N=10,000. The transmission rate is beta=0.3/day and the recovery rate is gamma=0.1/day (10-day infectious period). Initially 10 individuals are infected, with the rest susceptible. Questions:

- What is the basic reproduction number R0?
- When does the infection peak, and how many are infected at the peak?
- What fraction of the population is ultimately infected?
- What vaccination rate prevents an epidemic (herd immunity threshold)?

% ── Mapping ───────────────────────────────────────────────────────────
\subsection{Real-World Mapping}

\begin{longtable}{ll}
\toprule
\textbf{Real-World Concept} & \textbf{Mathematical Object} \\
\midrule
\endhead
healthy population & $S(t) — susceptible compartment$ \\
sick population & $I(t) — infected compartment$ \\
immune population & $R(t) — recovered compartment$ \\
contagiousness & $beta — transmission rate$ \\
recovery speed & $gamma — recovery rate$ \\
\bottomrule
\end{longtable}

% ── Constraints ───────────────────────────────────────────────────────
\subsection{Constraints}

\begin{enumerate}
  \item $\frac{dS}{dt} = -\frac{\beta S I}{N}$
  \hfill \textit{(susceptible individuals become infected)}
  \item $\frac{dI}{dt} = \frac{\beta S I}{N} - \gamma I$
  \hfill \textit{(infected gain from transmission, lose from recovery)}
  \item $\frac{dR}{dt} = \gamma I$
  \hfill \textit{(recovered individuals accumulate)}
  \item $S(t) + I(t) + R(t) = N$
  \hfill \textit{(population conservation)}
  \item $R_0 = \frac{\beta}{\gamma}$
  \hfill \textit{(basic reproduction number)}
  \item $v^* = 1 - \frac{1}{R_0}$
  \hfill \textit{(herd immunity threshold)}
\end{enumerate}

% ── Objective / Claim ─────────────────────────────────────────────────
\subsection{Objective}

\begin{equation}
\text{Solve } \frac{d\mathbf{y}}{dt} = f(t, \mathbf{y}), \quad \mathbf{y}(0) = \mathbf{y}_0
\end{equation}

% ── Approach ──────────────────────────────────────────────────────────
\subsection{Approach}

\begin{description}
  \item[Suggested approach:] Runge-Kutta 4(5) adaptive ODE integration (scipy.integrate.solve\_ivp)
  \item[Complexity class:] 
  \item[Available tools:] Runge-Kutta 4
\end{description}
% ── Phase B: Solution ────────────────────────────────────────────────
\section{Solution}

\begin{polyaresult}

\textbf{Answer}: R\_0 = 3.00; peak infected = 3,008; herd immunity = 66.7\%


\textbf{Optimal}: No / Unknown \hfill
\textbf{Feasible}: Yes
\textbf{Algorithm}: Runge-Kutta 4(5) adaptive ODE integration (scipy.integrate.solve\_ivp) \quad $$

\textbf{Time}: 0.13489555899286643s

\textbf{Certificate}: Population conservation verified. S(t)+I(t)+R(t)=N at all time steps.

\end{polyaresult}

% ── Solution Details ──────────────────────────────────────────────────
\subsection{Solution Details}

Population: N = 10,000
Transmission rate: beta = 0.3/day
Recovery rate: gamma = 0.1/day
Initial infected: I(0) = 10

Basic reproduction number: R\_0 = beta/gamma = 3.00
Peak infection: I\_max = 3,008 at t = 38.3 days
Final epidemic size (numerical): 94.1\% of population
Final epidemic size (theory): 94.0\% of population
Herd immunity threshold: v* = 1 - 1/R\_0 = 66.7\%

Vaccination scenarios:
  v = 0\%: R\_eff = 3.00, peak = 3,015, final size = 94.2\% [EPIDEMIC]
  v = 30\%: R\_eff = 2.10, peak = 1,204, final size = 57.7\% [EPIDEMIC]
  v = 50\%: R\_eff = 1.50, peak = 325, final size = 28.8\% [EPIDEMIC]
  v = 67\%: R\_eff = 0.99, peak = 10, final size = 1.5\% [contained]
  v = 80\%: R\_eff = 0.60, peak = 10, final size = 0.2\% [contained]

Disease-free equilibrium: unstable
Eigenvalues: 0.0, 0.0, 0.19999999999999998

% ── Verification ──────────────────────────────────────────────────────
\subsection{Verification}

\begin{longtable}{lcl}
\toprule
\textbf{Check} & \textbf{Status} & \textbf{Detail} \\
\midrule
\endhead
Population Conserved & \passmark & Passed \\
R0 Correct & \passmark & Passed \\
Peak Occurs & \passmark & Passed \\
Peak Time Reasonable & \passmark & Passed \\
Final Size Matches Theory & \passmark & Passed \\
Herd Immunity Correct & \passmark & Passed \\
Vaccination Prevents Epidemic & \passmark & Passed \\
All Compartments Nonneg & \passmark & Passed \\
\bottomrule
\end{longtable}

% ── Phase C: Interpretation ──────────────────────────────────────────
\section{Interpretation}

% ── The Question & The Answer ─────────────────────────────────────────
\subsection{The Question}
A disease outbreak in a population of N=10,000.

\subsection{The Answer}

\begin{polyainsight}[title={Bottom Line}]
R\_0 = 3.00; peak = 3,008 infected at day 38; herd immunity at 67\% vaccination.
\end{polyainsight}

% ── What This Means ───────────────────────────────────────────────────
\subsection{What This Means}

- R0: 3.0
- Peak infection: around day \textasciitilde{}30, with \textasciitilde{}2,900 infected simultaneously
- Final epidemic size: \textasciitilde{}94\% of the population ultimately infected
- Herd immunity threshold: 67\% vaccination rate prevents an epidemic
- Vaccination analysis: effective R0 drops below 1.0 at 67\% coverage, preventing large outbreaks

% ── Sensitivity Analysis ──────────────────────────────────────────────

% ── Figures ───────────────────────────────────────────────────────────

% ── Recommendations ───────────────────────────────────────────────────
\subsection{Recommendations}

\begin{enumerate}
  \item Monitor R\_0 closely; small changes in transmission can shift epidemic dynamics dramatically.
\end{enumerate}

% ── Limitations ───────────────────────────────────────────────────────
\subsection{Limitations}

\begin{polyawarnbox}
\begin{itemize}
  \item SIR model assumes homogeneous mixing; real populations have network structure.
  \item Parameters (beta, gamma) are assumed constant; real epidemics evolve over time.
\end{itemize}
\end{polyawarnbox}

% ── Appendix ─────────────────────────────────────────────────────────

\end{document}